\documentclass[12pt, a4paper, tocpage=plain]{abnt} % Fonte tamanho 12, papel A4, páginas do sumário sem o p.<número da página>

\usepackage[brazilian]{babel} % Gera datas e nomes em português com estilo brasileiro
\usepackage{hyperref} % Permite a criação de hyperlink no documento, como os links usados na referência
\usepackage[utf8]{inputenc} % Dá suporte para caracteres especiais como acentos e cedilha
\usepackage[T1]{fontenc}
\usepackage[alf]{abntcite} % Define o estilo de referência bibliográfica
\usepackage{graphicx} % Permite a utilização de imagens no documento
\usepackage[small]{caption} % Define as legendas das figuras com fontes menores do que o texto
\usepackage{pslatex} % Define que o formato da letra será Times New Roman
\usepackage{epigraph} % Permite a criação de epígrafes 
\usepackage{setspace} % Permite a definição de espaçamento entre linhas
\usepackage[top=3cm, left=3cm, right=2cm, bottom=2cm]{geometry} % Define as margens da folha

\setcounter{secnumdepth}{3} % Até três subsubsections numeradas
\setcounter{tocdepth}{3} % Até trẽs subsubsections numeradas

\setlength{\parindent}{1.25cm} % Define o recuo da primeira linha dos parágrafos para 1.25 cm

\usepackage{listings}
\usepackage{color}

\definecolor{dkgreen}{rgb}{0,0.6,0}
\definecolor{gray}{rgb}{0.5,0.5,0.5}
\definecolor{mauve}{rgb}{0.58,0,0.82}

\lstset{
  language=Python,% the language of the code
  numbers=left, %numeração de linhas à esquerda
  stepnumber=1,
  firstnumber=1,
  numberstyle=\tiny,
  extendedchars=true,
  frame=none,
  basicstyle=\footnotesize,
  stringstyle=\ttfamily,
  showstringspaces=false,
  captionpos=b,
  %language=Java, %deve ser definida na inclusão de cada trecho de código, 
  % pois podem existir linguagens diferentes em exemplos diferentes
  breaklines=true,
  breakautoindent=true,
  %estilos de comentário de uma e várias linhas
  keywordstyle=\color{blue},          % keyword style
  commentstyle=\color{dkgreen},       % comment style
  stringstyle=\color{mauve},         % string literal style
  frame=single,                   % adds a frame around the code
  extendedchars=\true,
  aboveskip=12pt,
  inputencoding=utf8,
}
\renewcommand{\lstlistingname}{Código}

\renewcommand{\ABNTchapterfont}{\bfseries} % Define a fonte do \chapter
\renewcommand{\ABNTchaptersize}{\large} % Define o tamanho da fonte do \chapter
\renewcommand{\ABNTsectionfontsize}{\large} % Define o tamanho da fonte da \section
\renewcommand{\ABNTsubsectionfontsize}{\large} % Define o tamanho da fonte do \subsection
\renewcommand{\ABNTsubsubsectionfontsize}{\large} % Define o tamanho da fonte do \subsubsection
\renewcommand{\ABNTbibliographyname}{REFERÊNCIAS BIBLIOGRÁFICAS} % Modifica o título gerado pelo \bibliographys

\begin{document} % Começo do TCC
\begin{titlepage}
 \begin{figure}[ht]
 \centering
 \scalebox{0.35}{\includegraphics{figuras/logo}}
 \end{figure}
 \begin{center}
   {\large BACHARELADO EM SISTEMAS DE INFORMAÇÃO} \\ [3.5cm]
   {\large COBIT: DS - ENTRAGA E SUPORTE } \\[0.6cm]
   {\large DS-5 - GARANTIR A SEGURANÇA DOS SISTEMAS }\\[9.0cm]
   \begin{flushright}
	{\large ALEXANDRE DOS SANTOS SAMPAIO} \\ [0.5cm]
	\begin{center}
		\begin{flushright}
			\textit{8º PERÍODO}
		\end{flushright}
	\end{center}
   \end{flushright}
   \vfill
   {\large Campos dos Goytacazes/RJ} \\
   {\large 2013}
 \end{center}
\end{titlepage} % Cria a capa
\begin{center}
\textbf{RESUMO}
\end{center}
\singlespacing

\noindent A Empresa \textbf{OFFXPRESS}, tem como objetivo a criação de um sistema de televendas para 80 pessoas.
Televenda refere-se a um conjunto de estratégias de divulgação e vendas de produtos e serviços pelo telefone. 
Através do telefone ocorre a apresentação de um serviço ou produto para um cliente, buscando a venda. Para tanto, 
como essa empresa não é especializada em TI e não possui o interesse em criar um departamento de TI, para tanto a 
empresa \textbf{Argentina Benzocriol Soluções em T.I.} , está sendo contratada para prestar todos os serviços de TI 
necessários.

\noindent Esta empresa deverá, portanto ser controlada durante todo o tempo de trabalho e para tanto serão definidos
os objetivos de controle sobre os quais esta cobrança será feita. 

\noindent Sendo assim, temos que criar uma estrutura para \textbf{\textit{governança da T.I}}. 

\noindent Neste momento iremos criar os objetivos de controle de Aquisição e Implementação, 
bem como os objetivos de controle de processos e de aplicativos. % Resumo do trabalho
\listoffigures % Gera o índice de figuras
\renewcommand{\contentsname}{SUMÁRIO} % Modifica o nome do sumário
\tableofcontents % Gera o sumário
\onehalfspacing % Define o espaçamento de 1.5cm entre linhas

\chapter{GOVERNANÇA DE TI}
\thispagestyle{empty}

\section{Requisitos de negócio para aparelhar o televendas}

\noindent Requisitos de negócio servem para explicar em termos do negócio 
	  o que deve ser entregue para agregar valor ao negócio propriamente 
	  dito.

\section{Objetivos de Negócio}

\noindent \subitem Criação de um canal de comunicação para contato direto entre o cliente e o atendente;
	  \subitem Uniformização e sistematização no atendimento;
	  \subitem Atuação em televendas ativa e receptiva;
	  \subitem Módulo para monitoramento de desempenho;
	  \subitem Criação de históricos de contatos e desempenho nas vendas;


\noindent

\end{document} % Fim do TCC
